\documentclass[11pt]{article}
\usepackage[a4paper,margin=1in]{geometry}
\usepackage[T1]{fontenc}
\usepackage{lmodern,microtype}
\usepackage{amsmath,amssymb,amsthm,mathtools}
\usepackage{booktabs,enumitem}
\usepackage{xcolor}
\usepackage[numbers,sort&compress]{natbib}
\usepackage[colorlinks=true,linkcolor=blue!55!black,citecolor=blue!55!black]{hyperref}

\newtheorem{theorem}{Theorem}[section]
\newtheorem{proposition}[theorem]{Proposition}
\newtheorem{corollary}[theorem]{Corollary}
\theoremstyle{definition}
\newtheorem{definition}[theorem]{Definition}
\theoremstyle{remark}
\newtheorem{remark}[theorem]{Remark}

\title{An Updated Deterministic-Envelope/Quotient Certificate Ledger}
\author{Bin Wang}
\date{July 2026}

\begin{document}
\maketitle

\begin{abstract}
We consolidate RH-262--RH-269 into a logically explicit certificate ledger.
The deterministic target now has a certified boundary constant, an exact
all-order parity dictionary, a direct order-29 logarithmic tail below
$2.6624745\times10^{-5}$, a unified envelope
$|a_n|<48q_*^n$, and the sharp law $a_n/q_*^n\to1$.  In contrast, the
archive supplies no legal anchored head, cloud coefficient bridge, or uniform
quotient tail; the sufficient continuum quotient criterion has none of its
four hypotheses verified.  Thus the five-component vector is
$(0,0,0,1,1)$, with exactly two satisfied obligations and zero complete
certificates in the audited route.  The zero count is scoped bookkeeping, not
a global nonexistence theorem.
\end{abstract}

\section{Objects and certificate semantics}

Let the Hardy-scaled deterministic target be normalized by
\begin{equation}
 \log G_H(z)=-\sum_{n\ge2}\frac{a_n}{n}z^n,
 \qquad r_H=\frac{17}{20}.
 \label{eq:target}
\end{equation}
For a selected moving-cloud or quotient family, write its candidate
coefficients as $\tau_{n,s}$.  A complete head--tail certificate requires
more than a target estimate: it must identify a legal finite head, prove that
its coefficients match the deterministic target in the stated protocol, and
control both quotient and target tails uniformly.

\begin{definition}[Five-obligation ledger]\label{def:ledger}
Use the fixed order $(H,B,Q,A,M)$, where
\begin{equation}
\begin{aligned}
 H&=\text{legal anchored head},&
 B&=\text{coefficient bridge},\\
 Q&=\text{uniform quotient tail},&
 A&=\text{analytic target tail},\\
 M&=\text{certified target boundary constant}.&&
\end{aligned}
\label{eq:vector}
\end{equation}
the complete-certificate predicate is the conjunction
$H\wedge B\wedge Q\wedge A\wedge M$.
\end{definition}

The five entries are deliberately independent.  In particular, a theorem
about the deterministic sequence $a_n$ cannot by itself establish that a
noise-dependent spectral selection produces those coefficients.  This is the
same firewall that separated the deterministic target from the moving-cloud
trace-envelope question in RH-241 \citep{WangRH241}.

\section{The deterministic target stack}

The target side now consists of exact analytic identities and certified
interval bounds rather than a finite root-rate extrapolation.  RH-262 proves,
using the RH-15 factorization and the RH-13 Arb certificate
\citep{WangRH13,WangRH15,WangRH262},
\begin{equation}
 M_{7/5}:=\sup_{|z|=7/5}|\log G_H(z)|<107.906078<108.
 \label{eq:boundary}
\end{equation}
RH-263 then gives the all-order parity dictionary \citep{WangRH263}:
for odd $n\ge3$,
\begin{equation}
 a_n=\frac{(r_H\lambda)^{-n}}{1+\lambda^{-n}},
 \label{eq:odd}
\end{equation}
while for $k\ge1$,
\begin{equation}
 a_{2k}=r_H^{-2k}\left[
 2\operatorname{tr}(T^k)+\frac{2\lambda^{-2k}}{1+\lambda^{-k}}
 -\frac{\lambda^{-2k}}{1-\lambda^{-2k}}
 \right].
 \label{eq:even}
\end{equation}
The comparison against the 27 archived RH-253 rows through order 28 has
maximum floating residual $6.442918843840156\times10^{-14}$; the proof of
\eqref{eq:odd}--\eqref{eq:even} is factorwise and does not depend on that
finite check \citep{WangRH253}.

Put
\begin{equation}
 q_*=(r_H\lambda)^{-1}=0.7008752258547757\ldots,
 \qquad \rho_*=q_*^{-1}=1.4267874838640739\ldots .
 \label{eq:qstar}
\end{equation}
RH-267 proves the unified all-order envelope
\begin{equation}
 |a_n|<48q_*^n\qquad(n\ge2),
 \label{eq:envelope}
\end{equation}
and RH-268 proves its base is sharp in the deterministic sense
\citep{WangRH267,WangRH268}:
\begin{equation}
 \frac{a_n}{q_*^n}\longrightarrow1.
 \label{eq:sharp}
\end{equation}
Thus no bound $|a_n|\le Cq^n$ can hold for every $n$ with $q<q_*$, although
the prefactor $48$ is not claimed to be optimal.

\begin{proposition}[Envelope-only logarithmic tail]\label{prop:envtail}
If $0\le R<\rho_*$ and $N\ge2$, then \eqref{eq:envelope} implies
\begin{equation}
 \sum_{n\ge N}\frac{|a_n|R^n}{n}
 <\frac{48(q_*R)^N}{N(1-q_*R)}.
 \label{eq:envtail}
\end{equation}
At $R=1$ and $N=29$, the certified comparison $q_*<0.700876$ gives the
clean bound
\begin{equation}
 \sum_{n\ge29}\frac{|a_n|}{n}<0.000184751.
 \label{eq:coarsetail}
\end{equation}
\end{proposition}

\begin{proof}
For $n\ge N$, replace $1/n$ by $1/N$ in the positive majorant and sum the
geometric series with ratio $q_*R<1$.  Substituting the outward safe value
$0.700876$ gives \eqref{eq:coarsetail}.
\end{proof}

The direct factorwise certificate of RH-264 is stronger at the operational
order \citep{WangRH264}:
\begin{equation}
 \sum_{n\ge29}\frac{|a_n|}{n}<0.000026624745,
 \qquad e^{\sum_{n\ge29}|a_n|/n}-1<0.000026625100.
 \label{eq:directtail}
\end{equation}
Its reported logarithmic endpoint is more than $6.93$ times smaller than the
coarse envelope-only endpoint in \eqref{eq:coarsetail}.  RH-265 extends this
factorwise construction to a certified tail ladder; only the $N=29$ row is
aligned with the currently archived order-28 head, and every row $N\ge37$
remains a conditional interface rather than a constructed head
\citep{WangRH265}.

\begin{theorem}[Deterministic target package]\label{thm:target}
The archived target has a certified analytic tail and certified boundary
constant.  Its coefficient base is exactly $q_*$, its logarithmic radius is
exactly $\rho_*$, and its current order-29 tail obeys \eqref{eq:directtail}.
These conclusions concern the deterministic sequence $a_n$ only and do not
imply a moving-cloud coefficient bridge.
\end{theorem}

\begin{proof}
The boundary and direct tail are \eqref{eq:boundary} and
\eqref{eq:directtail}.  Equations \eqref{eq:envelope}--\eqref{eq:sharp} give
the exact root rate and, by Cauchy--Hadamard, the radius $\rho_*$.  None of
these statements contains an identification $\tau_{n,s}=a_n$ or a limiting
version of that identity, so the bridge does not follow.
\end{proof}

\section{The quotient side remains conditional}

RH-259 has 23 finite quotient blocks: every twelfth power is contractive,
none of the one-step blocks is contractive, and nine archived endpoints are
missing \citep{WangRH259}.  RH-266 proves the exact logical limitation of
this evidence: finitely many pointwise contractions do not imply a continuum
uniform bound without a modulus or an interval-family enclosure
\citep{WangRH266}.  This does not prove that the actual quotient family is
nonuniform.

RH-269 supplies a sufficient continuum theorem \citep{WangRH246,WangRH269}.
It requires
\begin{enumerate}[label=(\roman*)]
 \item Hilbert--Schmidt convergence of the noisy operators;
 \item a common finite-rank contour isolating the spectral cluster;
 \item a uniform resolvent bound on that contour;
 \item a contractive power of the limiting orthogonal quotient.
\end{enumerate}
Together these hypotheses stabilize the Riesz and orthogonal projections and
produce uniform RH-246 constants $K_m$, $\eta_m<1$, and $L_r$.  The current
archive verifies none of the four hypotheses as a continuum theorem.  Hence
the criterion is exact but inactive, and the uniform quotient tail remains
open.

\section{Exact updated ledger}

\begin{center}
\small
\begin{tabular}{p{0.46\linewidth}p{0.39\linewidth}}
\toprule
component & archived status after RH-269 \\
\midrule
certified target boundary constant & true; \eqref{eq:boundary} \\
exact deterministic parity anchor & true; \eqref{eq:odd}--\eqref{eq:even} \\
direct order-29 deterministic tail & true; \eqref{eq:directtail} \\
deterministic all-order envelope & true; \eqref{eq:envelope} \\
sharp deterministic base and radius & true; \eqref{eq:sharp} \\
legal anchored head & not supplied in audited classes \\
cloud coefficient bridge & open \\
uniform quotient tail & open \\
RH-269 criterion hypotheses & $0/4$ verified \\
\bottomrule
\end{tabular}
\end{center}

\begin{theorem}[Updated complete-certificate ledger]\label{thm:ledger}
For Definition~\ref{def:ledger}, the archived truth vector is
\begin{equation}
 \boxed{(H,B,Q,A,M)=(0,0,0,1,1).}
 \label{eq:status}
\end{equation}
Exactly two of five obligations are satisfied.  The complete certificate
count in the audited route is therefore zero.
\end{theorem}

\begin{proof}
RH-262 certifies $A=M=1$.  RH-263--RH-268 strengthen the target side but do
not change the three cloud/quotient entries.  The archived head ledger still
contains no passing legal anchored head, so $H=0$; no cloud-to-target
coefficient theorem gives $B=1$; and RH-266--RH-269 leave $Q=0$, with the
sufficient criterion at $0/4$ hypotheses.  The complete predicate is the
conjunction of all five entries, so \eqref{eq:status} is incomplete and its
audited complete-certificate count is zero.
\end{proof}

\begin{remark}[Scope of the zero count]
The theorem does not exclude a future selector outside the audited head
classes, an operator-derived non-idempotent quotient, a cloud bridge, or a
continuum proof of the RH-269 hypotheses.  It says exactly that no complete
certificate is assembled from the presently archived components.
\end{remark}

\section{Protocol and claim boundary}

The companion audit reads the eight structured result files RH-262--RH-269,
checks their status fields, theorem boundaries, numerical endpoints, and all
forty inherited Gate A--E flags, and writes a deterministic ledger.  It does
not rerun long operator calculations, refit finite coefficients, or infer
uniformity between samples.

The deterministic all-order envelope and sharp base are not moving-cloud
theorems.  The 27-row parity cross-check is not used as proof of an all-order
cloud identity, and the 23 finite contractions are not a uniform quotient
certificate.  Gates A--E remain false/open.  No Hilbert--Polya operator is
constructed, no Riemann zeros are identified, no equality with the completed
zeta divisor is asserted, and no implication of the Riemann Hypothesis is
claimed.

\bibliographystyle{plainnat}
\bibliography{references}
\end{document}
